\subsection{CUDA ASR Results}
\label{sec:cuda_results}

All four custom CUDA task implementations passed their
respective correctness tests. However, numerical
correctness did not imply end-to-end speedup.
Table~\ref{tab:asr_ablation} reports controlled
comparisons on the H200 MIG partition.

\begin{table}[!htbp]
\centering
\caption{End-to-end ASR performance on H200 MIG.
Latency is the mean of three timed runs after one
warm-up. All configurations used the same input.}
\label{tab:asr_ablation}
\small
\begin{tabular}{lrr}
\hline
Configuration & Latency (ms) & Peak GiB \\
\hline
PyTorch baseline & 1339.2 & 8.90 \\
All custom CUDA & 5867.1 & 8.56 \\
PyTorch GEMM & 2015.5 & 8.59 \\
PyTorch GEMM + attention & \textbf{1296.2} & 8.90 \\
\hline
\end{tabular}
\end{table}

The fully custom implementation was approximately
$4.38\times$ slower than the PyTorch baseline,
although its measured peak GPU memory was
0.34~GiB lower. Replacing custom GEMM with
PyTorch's implementation reduced latency by
3851.6~ms (65.6\% relative to the fully custom
configuration). Replacing custom attention as
well saved a further 719.3~ms (35.7\% relative
to the preceding hybrid).

The resulting hybrid was 3.2\% faster than the
PyTorch-only baseline on this input. This small
difference should not be interpreted as a
general speedup: the experiment used one
transcription workload and three timed runs.

\paragraph{Profiling and bottleneck attribution.}
Profiling the fully custom configuration attributed
4966.6~ms (84.5\%) to linear operations across
2757 calls, and 882.0~ms (15.0\%) to attention
across 396 calls. These results agree with the
controlled ablations: GEMM was the principal
performance bottleneck, followed by attention.
Kernel-launch overhead and limited reuse in
the simple tiled implementation are plausible
contributors, while the measurements alone do
not isolate their individual effects.

\paragraph{Tile-size sensitivity.}
For the encoder \texttt{fc1} workload, increasing
GEMM tile size from 8 to 32 reduced latency
from 32.459 to 18.389~ms, a $1.77\times$
improvement. In contrast, for \texttt{lm\_head}
with $M=1$, tile 8 took 2.911~ms while tile 32
took 7.141~ms. Larger square tiles waste
threads when the output has only one row,
illustrating why the optimal tile size depends
on matrix shape. The final submitted GEMM
implementation uses tile size 16.

These observations motivate specialised
implementations for different matrix shapes
rather than relying on one fixed tile size.
